\documentclass[11pt]{article}
\usepackage[margin=1.05in]{geometry}
\usepackage{amsmath,amssymb,amsthm,mathtools}
\usepackage{booktabs,microtype,xcolor,graphicx}
\usepackage[shortlabels]{enumitem}
\usepackage[colorlinks=true,linkcolor=blue!45!black,citecolor=blue!45!black,urlcolor=blue!45!black]{hyperref}
\hypersetup{pdftitle={How much does the endpoint say about the switching rate?},pdfauthor={Arjun Pemmasani}}
\newtheorem{theorem}{Theorem}[section]
\newtheorem{lemma}[theorem]{Lemma}
\newtheorem{proposition}[theorem]{Proposition}
\newtheorem{corollary}[theorem]{Corollary}
\theoremstyle{definition}
\newtheorem{definition}[theorem]{Definition}
\newtheorem{example}[theorem]{Example}
\theoremstyle{remark}
\newtheorem{remark}[theorem]{Remark}
\newcommand{\Z}{\mathbb Z}
\newcommand{\E}{\mathbb E}
\newcommand{\Prob}{\mathbb P}
\newcommand{\Var}{\operatorname{Var}}
\newcommand{\Cov}{\operatorname{Cov}}
\newcommand{\ones}{\mathbf 1}
\emergencystretch=1.5em
\title{How much does the endpoint say about the switching rate?\\[4pt]
\large Working note for direction S1 (moved out of Paper~C, October 2026)}
\author{}
\date{October 2026}
\begin{document}
\maketitle

\section{Setting}\label{sec:setting}

This note was Section~8 and Appendix~E of an earlier version of Paper~C
(``Face selection for slowly switching Markov chains and the planar persistent
walk''). The mathematics is unchanged. It uses only Paper~A~\cite{paperA},
and this section recalls what it needs from there. Smiga, Radaelli, Binder and
Landi~\cite{smiga2023} compute the Fisher information of the empirical
distribution under a Gaussian approximation. Their leading term vanishes for
the symmetric walk of Paper~A, whose stationary law does not depend on the
switching rate.

Kagey's board~\cite{kagey131} has $N$ steps. The first step is left or right
with probability $\frac12$, and each later step repeats the previous one with
probability $p$. We write $q=1-p$ and $u=q/p$, the odds of a switch. We write
$L=\log N$, and $T=Nq$ is the expected number of switches (up to the first
step). $P_N(k)$ is the probability of $k$ right steps, that is, that the ball
lands in bin $k$. The endpoints have $P_N(0)=P_N(N)=p^{N-1}/2$.

For $a,b\ge1$ let $S_{a,b}(u)=\sum_wu^{s(w)}$, the sum over the words $w$ with
$a$ right steps and $b$ left steps, where $s(w)$ is the number of switches of
$w$. Then $S_{a,b}(q/p)$ is the ratio $P_N(a)/P_N(0)$ of the probabilities of
bin $a$ and bin $0$ after $a+b=N$ steps (the run count of~\cite{paperA}). Let
$S_N=S_{\lfloor N/2\rfloor,\lceil N/2\rceil}$. It has nonnegative coefficients
and no constant term, so $S_N(u)=1$ has exactly one positive root $u_N$, and
$p_N=1/(1+u_N)$ is the crossing of the middle and the endpoints. Put
$z_N=N(1-p_N)$ and $c_0=\frac12\log(\pi/8)$. Paper~A proves that
$z_N=L-\frac12\log L+c_0+O(\log L/L)$. We call this the second-order expansion
of $z_N$.

\section{The information in the endpoint}\label{sec:fisher}

This note asks a statistical question about Kagey's board. If we only see the
bin where the ball lands, how much do we learn
about the switching probability? The answer is that the bin keeps about a
fraction $1/(2T)$ of the Fisher information of the whole path, so about
$1/(2\log N)$ at Kagey's crossing. We use the notation of
Section~\ref{sec:setting}. The first step is left or right with probability
$\frac12$ each, and each later step repeats the one before it with probability
$p=1-q$, where $0<q<1$, and $T=Nq$.
Put $\rho=1-2q$, $u=q/p$ and $y=Nu=T/p$. Let $J$ be the number of
switches, $K$ the number of right steps, $X=2K-N$ the endpoint and $W=X/N$. A path with $J$ switches has probability
$\frac12p^{N-1-J}q^J$. So $J\sim\operatorname{Bin}(N-1,q)$, and the path
carries the Fisher information $(N-1)/(pq)$ about $q$. The endpoint carries
$I_K(q)=\sum_{k=0}^N(\partial_q\Prob(K=k))^2/\Prob(K=k)$, and we call
$e_N(q)=I_K(q)\,pq/(N-1)$ the \emph{efficiency} of the endpoint. It is the
fraction of the information left when we see only the endpoint.

The endpoint sees the switches only through how far the walk spreads. A
normal $X$ with mean $0$ and variance $Np/q$, close to $\Var X$ for large
$T$, would carry the information $\frac12(\partial_q\log(Np/q))^2=1/(2p^2q^2)$,
which is the fraction $1/(2(N-1)pq)\approx1/(2T)$. Theorem~\ref{thm:fisher}
confirms this.

\begin{proposition}[a correlation ratio]\label{prop:eta}
For $N\ge2$ and $0<q<1$,
\[
 e_N(q)=\frac{\Var(\E[J\mid K])}{\Var J}.
\]
For $0<k<N$ we have $\E[J\mid K=k]=\frac{d}{dt}\log S_{k,N-k}(e^t)$ at
$t=\log(q/p)$, where $S_{a,b}$ is the polynomial of
Section~\ref{sec:setting}. For $k\in\{0,N\}$ we have $J=0$.
\end{proposition}

\begin{proof}
The score of a path is $(J-(N-1)q)/(pq)$. Since $\Prob(K=k)$ is a sum of path
probabilities, $\partial_q\log\Prob(K=k)=\E[(J-(N-1)q)/(pq)\mid K=k]$. As
$\E J=(N-1)q$, this gives $I_K(q)=\Var(\E[J\mid K])/(pq)^2$, and we divide by
$(N-1)/(pq)=\Var J/(pq)^2$. By the run count of~\cite{paperA}, for $0<k<N$
the probability $\Prob(K=k,J=j)$ is $\frac12p^{N-1}u^j$ times the number of
words with $k$ right steps, $N-k$ left steps and $j$ switches. Summing over
$j$ gives $\Prob(K=k)=\frac12p^{N-1}S_{k,N-k}(u)$, so
$\E[J\mid K=k]=uS'_{k,N-k}(u)/S_{k,N-k}(u)$. Only the two constant paths have
$K\in\{0,N\}$.
\end{proof}

This identity, checked in Lean, says that $e_N$ is the correlation ratio of
$J$ on $K$. Since $\E[J\mid K]$ is the orthogonal projection of $J$ onto
the functions of $K$ in $L^2$, the Cauchy--Schwarz inequality gives
\begin{equation}\label{eq:fisher-proj}
 e_N(q)\ge\frac{\Cov(J,g(K))^2}{\Var J\,\Var g(K)}
 \qquad\text{whenever }\Var g(K)>0.
\end{equation}

\begin{proposition}[few switches]\label{rem:fisher-small}
\begin{enumerate}[(a)]
\item For fixed $N\ge3$, $e_N(q)=1-\frac{N-2}2q+O_N(q^2)$ as $q\to0$.
\item Let $q\to0$ with $T=Nq$ bounded. Then $e_N(q)\to1$ if and only if
$T\to0$.
\end{enumerate}
\end{proposition}

In words, the endpoint loses information as soon as paths with two switches
appear. The reason is that one switch and two switches both leave the endpoint
uniform on the interior bins, so an interior bin cannot tell them apart.

\begin{proof}
By Proposition~\ref{prop:eta}, $1-e_N=\E\Var(J\mid K)/\Var J$, and
$\Var J=(N-1)pq$. Fix an interior bin $0<k<N$. Exactly two words with $k$
right steps have one switch, and exactly $N-2$ have two switches ($k-1$ that
start to the right and $N-k-1$ that start to the left). So
\[
 \Prob(J=1,K=k)=p^{N-2}q,\qquad\Prob(J=2,K=k)=\frac{N-2}2\,p^{N-3}q^2 .
\]

(a) Here $N$ is fixed. Given $K=k$ the number of switches is $1$, or $2$ with
probability $\frac{N-2}2q+O_N(q^2)$, or larger with probability $O_N(q^2)$. So
$\Var(J\mid K=k)=\frac{N-2}2q+O_N(q^2)$, and
$\Prob(0<K<N)=(N-1)q+O_N(q^2)$. For $K\in\{0,N\}$ the conditional variance is
$0$. Hence $\E\Var(J\mid K)=\frac{(N-1)(N-2)}2q^2+O_N(q^3)$, and dividing by
$(N-1)pq$ gives the claim.

(b) For the upper bound we replace $\E[J\mid K]$ by the indicator of an
interior bin, which can only raise the mean squared error. This gives
\[
 \E\Var(J\mid K)\le\E\bigl(J-\mathbf 1\{0<K<N\}\bigr)^2
 =\E\bigl[(J-1)^2;J\ge2\bigr]\le\E[J(J-1)]=(N-1)(N-2)q^2 ,
\]
because $J\ge1$ exactly on the interior bins. So
$1-e_N\le(N-2)q/p\le T/p$, which tends to $0$ if $T\to0$.

For the lower bound, a variable that takes two values $1$ apart with
probabilities $a$ and $b$ has variance at least $ab$. So
$\Var(J\mid K=k)\ge\Prob(J=1\mid K=k)\Prob(J=2\mid K=k)$, and by the
Cauchy--Schwarz inequality
\[
 \E\Var(J\mid K)\ge\sum_{0<k<N}\frac{\Prob(J=1,K=k)\Prob(J=2,K=k)}{\Prob(K=k)}
 \ge\Bigl(\sum_{0<k<N}\sqrt{\Prob(J=1,K=k)\Prob(J=2,K=k)}\Bigr)^2
 =\frac{(N-1)^2(N-2)}2\,p^{2N-5}q^3 .
\]
Dividing by $(N-1)pq$ gives $1-e_N\ge\frac{(N-1)(N-2)}2p^{2N-6}q^2$. If $T$
stays in an interval $[t_0,t_1]$ with $t_0>0$, the right side tends to a limit
between $\frac12t_0^2e^{-2t_1}$ and $\frac12t_1^2$, and it is positive. So
$e_N$ does not tend to $1$ along any sequence on which $T$ stays away from
$0$.
\end{proof}

\begin{theorem}[many switches]\label{thm:fisher}
Let $R^2_N(q)=\Cov(J,X^2)^2/(\Var J\,\Var X^2)$ be the squared correlation
of $J$ and $X^2$. For $N\ge2$ and $0<q<1$,
\[
 R^2_N(q)=\frac1{2(N-1)pq}\cdot
 \frac{\bigl[1+\rho^N-p(1-\rho^N)/T\bigr]^2}{V},
\]
where
\[
 V=1-\frac{1+8\rho+\rho^2+8\rho^{N+1}}{4pT}
 +\frac{\rho(1-\rho^N)\bigl[2(1+2\rho)(2+\rho)-\rho(1-\rho^N)\bigr]}{8p^2T^2}.
\]
There are absolute constants $C$ and $N_0$ such that for $N\ge N_0$,
$0<q\le\frac12$ and $1\le T\le N^{1/4}$,
\[
 0\le e_N(q)-R^2_N(q)\le C\Bigl(\frac1{T^3}+\frac{T^3}N\Bigr).
\]
Consequently, in the same range and with an absolute implied constant,
\[
 e_N(q)=\frac1{2T}+\frac1{4T^2}+O\Bigl(\frac1{T^3}+\frac{T^3}N\Bigr).
\]
\end{theorem}

In words, when $T$ is large but at most $N^{1/4}$, the endpoint keeps about the
fraction $1/(2T)$ of the information, and almost all of what it keeps is carried
by the single statistic $X^2$. The lower bound $e_N\ge R^2_N$ is
\eqref{eq:fisher-proj} with $g(K)=X^2$,
for all $N\ge2$ and $0<q<1$. For the upper bound, Proposition~\ref{prop:eta}
and the uniform Bessel estimate for every bin in~\cite{paperA}
(recalled in~\eqref{eq:fisher-bessel-bound} below) make
$\E[J\mid K]$ close to $y\sqrt{1-W^2}+\frac12$. As $W$ is of order
$T^{-1/2}$, this is close to $y(1-W^2/2)+\frac12$, which lies in the span of
$1$ and $X^2$. Appendix~\ref{app:fisher} proves the closed form and the
bound.

\begin{proof}[Proof of the expansion in Theorem~\ref{thm:fisher}]
If $T\le18$, then $\lvert e_N-\frac1{2T}-\frac1{4T^2}\rvert\le1\le18^3T^{-3}$,
since $e_N$ and $\frac1{2T}+\frac1{4T^2}$ lie in $[0,1]$. Let $T\ge18$. As $q\le\frac12$, we have $p\ge\frac12$
and $0\le\rho<1$, and $\rho^N\le e^{-2Nq}=e^{-2T}\le T^{-2}$. The middle term of $V$ lies
in $[0,9/T]$, and since $1+8\rho+\rho^2=10-20q+4q^2$ it equals
$\frac5{2T}+O(q/T+T^{-2})$. The last term lies in $[0,9/T^2]$. So
$V\ge\frac12$, $1/V=1+\frac5{2T}+O(q/T+T^{-2})$, the squared bracket is
$1-\frac2T+O(q/T+T^{-2})$, and $1/(2(N-1)pq)=(1+O(q+1/N))/(2T)$. Multiplying,
$R^2_N=\frac1{2T}+\frac1{4T^2}+O(T^{-3}+q/T+1/(NT))$ with absolute
constants. Since $q/T=1/N\le T^3/N$, the bound on $e_N-R^2_N$ gives the
expansion.
\end{proof}

\begin{corollary}[at Kagey's crossing]\label{cor:fisher-crossing}
Let $p_N$ be the value of $p$ at which the middle bin and the endpoints are
equally likely, and let $q_N=1-p_N$ and $z_N=Nq_N$
(Section~\ref{sec:setting}). Then
\[
 e_N(q_N)=\frac1{2L}+\frac{\log L+1+\log(8/\pi)}{4L^2}
 +O\Bigl(\frac{(\log L)^2}{L^3}\Bigr).
\]
\end{corollary}

\begin{proof}
Since $S_N(u)\ge2u$, we have $u_N\le\frac12$, so $p_N\ge\frac23$ and
$q_N\le\frac13$. By the second-order expansion of $z_N$ in~\cite{paperA},
$z_N=L(1-\varepsilon)$ with
$\varepsilon=(\frac12\log L-c_0)/L+O(\log L/L^2)$. So Theorem~\ref{thm:fisher}
applies with $T=z_N$ for large $N$ and gives
$e_N(q_N)=\frac1{2z_N}+\frac1{4z_N^2}+O(L^{-3})$. Here
$\frac1{2z_N}=\frac1{2L}+\frac{\log L-2c_0}{4L^2}+O((\log L)^2/L^3)$,
$\frac1{4z_N^2}=\frac1{4L^2}+O(\log L/L^3)$ and $-2c_0=\log(8/\pi)$.
\end{proof}

So $e_N\sim1/(2\log N)$ there, but the relative correction is of order
$(\log\log N)/\log N$ and is large at every $N$ we computed. The product
$2Le_N$ is $1.74$ at $N=100$ and $1.35$ at $N=3200$ (Table~\ref{tab:fisher}).

\begin{table}[ht]
\centering\small
\begin{tabular}{@{}rcccccc@{}}
\toprule
$N$ & $z_N$ & $e_N(q_N)$ & $2z_Ne_N$ & $2Le_N$ & $R^2_N(q_N)$
 & $\frac1{2z_N}+\frac1{4z_N^2}$\\
\midrule
50 & 2.8998 & 0.24117 & 1.3987 & 1.8869 & 0.20873 & 0.20216\\
100 & 3.5034 & 0.18869 & 1.3221 & 1.7379 & 0.16642 & 0.16309\\
400 & 4.7501 & 0.12681 & 1.2047 & 1.5196 & 0.11715 & 0.11634\\
1600 & 6.0246 & 0.09433 & 1.1366 & 1.3919 & 0.09007 & 0.08988\\
3200 & 6.6684 & 0.08361 & 1.1151 & 1.3496 & 0.08070 & 0.08060\\
\bottomrule
\end{tabular}
\caption{The efficiency of the endpoint at Kagey's crossing $q_N=1-p_N$,
computed from the exact law, with $z_N=Nq_N$ and $L=\log N$. The last two
columns are the squared correlation of Theorem~\ref{thm:fisher} and its
expansion to two terms.}
\label{tab:fisher}
\end{table}\subsection*{Many balls}

Now drop $n$ independent balls and record only the histogram of their bins.
How well can $q$ be estimated? By the Cram\'er--Rao bound no unbiased
estimator has variance below $1/(nI_K(q))$, and the maximum likelihood
estimator attains this as $n\to\infty$. The next corollary says that a much
simpler estimator does almost as well. Put $m(q)=\E_qX^2$, a polynomial in
$q$.

\begin{corollary}[the moment estimator]\label{cor:fisher-moment}
Let $X_1,\dots,X_n$ be the endpoints of $n$ independent balls with switching
probability $q$, and suppose that $R^2_N(q)>0$. Then $m'(q)\ne0$, and with
probability tending to $1$ the equation
$m(\hat q)=\frac1n\sum_{i=1}^nX_i^2$ has a solution $\hat q_n$ that converges
to $q$ in probability. For this solution
\[
 \sqrt n\,(\hat q_n-q)\Rightarrow N\Bigl(0,\frac{pq}{(N-1)R^2_N(q)}\Bigr)
 \qquad(n\to\infty).
\]
So its asymptotic efficiency, relative to the bound $1/(nI_K(q))$, is
$R^2_N(q)/e_N(q)$. There is an absolute constant $C'$ such that, with $N_0$
as in Theorem~\ref{thm:fisher}, for $N\ge N_0$, $0<q\le\frac12$ and
$1\le T\le N^{1/4}$,
\[
 1-C'\Bigl(\frac1{T^2}+\frac{T^4}N\Bigr)\le\frac{R^2_N(q)}{e_N(q)}\le1 .
\]
\end{corollary}

In words, among all estimators that see only the bins, matching the mean
square of the endpoint loses a fraction of order $1/T^2+T^4/N$ of the
information. At Kagey's crossing the efficiency $R^2_N/e_N$ is $0.87$ at
$N=50$ and $0.97$ at $N=3200$ (Table~\ref{tab:fisher}).

\begin{proof}
The law of a path is a polynomial in $q$ with score $(J-(N-1)q)/(pq)$, so
$m'(q)=\E[X^2(J-(N-1)q)]/(pq)=\Cov(J,X^2)/(pq)$. This is not zero because
$R^2_N(q)>0$. So $m$ is strictly monotone near $q$. By the law of large
numbers the mean of the $X_i^2$ tends to $m(q)$, so with probability tending
to $1$ it lies in the image of a small interval around $q$, and the solution
in that interval converges to $q$. The central limit theorem and the delta
method give the limit law with variance
\[
 \frac{\Var X^2}{m'(q)^2}=\frac{(pq)^2\Var X^2}{\Cov(J,X^2)^2}
 =\frac{(pq)^2}{R^2_N(q)\Var J}=\frac{pq}{(N-1)R^2_N(q)} .
\]
Since $I_K(q)=e_N(q)(N-1)/(pq)$, the ratio of $1/I_K(q)$ to this variance is
$R^2_N/e_N$, and it is at most $1$ by~\eqref{eq:fisher-proj}. For the lower
bound let $C_1$ be the constant in the expansion of
Theorem~\ref{thm:fisher}, and put $C'=4\max(C,C_1)$ and
$\epsilon=T^{-2}+T^4/N$. If $C'\epsilon\ge1$ the claim is trivial, since
$R^2_N/e_N\ge0$. Otherwise $C_1\epsilon<\frac14$, so the expansion gives
$e_N\ge\frac1{2T}-\frac{C_1\epsilon}T\ge\frac1{4T}$. Then
$1-R^2_N/e_N=(e_N-R^2_N)/e_N\le4T\cdot C\epsilon/T=4C\epsilon\le C'\epsilon$.
\end{proof}


\appendix
\section{Proof of the Fisher information theorem}\label{app:fisher}

We keep the notation of Section~\ref{sec:fisher}. Let
$\sigma_i\in\{-1,1\}$ be the direction of step $i$. Then
$X=\sigma_1+\dots+\sigma_N$, and $J$ counts the $i<N$ with
$\sigma_{i+1}\ne\sigma_i$. Put $Z=\E[J\mid K]$, so $e_N=\Var Z/\Var J$
by Proposition~\ref{prop:eta}.

\begin{lemma}[exact moments]\label{lem:fisher-moments}
For $N\ge1$ and $0<q<1$,
\begin{align*}
 \E X^2&=\frac{Np}q-\frac{\rho(1-\rho^N)}{2q^2},\\
 \Cov(J,X^2)&=-\frac p{q^2}\bigl[Nq(1+\rho^N)-p(1-\rho^N)\bigr],\\
 \E X^4&=\frac{3N^2p^2}{q^2}-\frac{Np(1+10\rho+\rho^2)}{2q^3}
 +\frac{\rho(1+2\rho)(2+\rho)(1-\rho^N)}{2q^4}
 -\frac{3Np\rho^{N+1}}{q^3},
\end{align*}
and $\Var X^2=2N^2p^2V/q^2$, with $V$ as in Theorem~\ref{thm:fisher}.
\end{lemma}

\begin{proof}
We show that each moment, as a sequence in $N$, is a combination of the
sequences $N^i$ and $N^i\rho^N$ with $i\le4$. Two such combinations that agree
at ten values of $N$ agree everywhere, and a script compares them at those
values exactly.

Let $m_N(\phi,\chi)=\E e^{\phi X+\chi J}$. With
$A_{ab}=\Prob(\sigma_2=b\mid\sigma_1=a)\,e^{\phi b+\chi\mathbf 1[a\ne b]}$ and
$v_a=\frac12e^{\phi a}$ we have $m_N=v^\top A^{N-1}\ones$. So
$\sum_{N\ge1}m_Nz^N=U/D$ as formal power series, where
$U(z)=zv^\top\operatorname{adj}(I-zA)\ones$ and $D(z)=\det(I-zA)$ have degree
at most $2$ in $z$. For a partial derivative
$\partial$ of order $m$ in $(\phi,\chi)$, induction and the quotient rule give
$\partial(U/D)=U_\partial/D^{m+1}$ with $\deg U_\partial\le2m+2$, as each
derivative multiplies the numerator by $D$ or a derivative of $D$. At
$\phi=\chi=0$ the matrix $A$ has trace $1+\rho$ and determinant $\rho$, so
$D(z)=(1-z)(1-\rho z)$. Let $q\ne\frac12$, so its roots $1$ and
$1/\rho$ are distinct. By partial
fractions $U_\partial/D^{m+1}$ is a constant plus a combination of
$(1-z)^{-j}$ and $(1-\rho z)^{-j}$ with $j\le m+1$, whose coefficients of
$z^N$ are $\binom{N+j-1}{j-1}$ and $\binom{N+j-1}{j-1}\rho^N$. So for $N\ge1$
the sequence $\partial m_N$ at $0$ lies in
$\mathcal S_m=\operatorname{span}\{N^i,N^i\rho^N:i\le m\}$. A sequence in
$\mathcal S_m$ satisfies the recurrence with characteristic polynomial
$(t-1)^{m+1}(t-\rho)^{m+1}$, so it is $0$ if it vanishes at
$N=1,\dots,2m+2$.

The moments $\E X^2$, $\E[JX^2]$ and $\E X^4$ are derivatives of $m_N$ of
order at most $4$ at $0$, so they lie in $\mathcal S_4$. The claimed
formulas for $\E X^2$, $\E X^4$ and $\E[JX^2]=\Cov(J,X^2)+(N-1)q\,\E X^2$
lie in $\mathcal S_2$. So it is enough to check them at $N=1,\dots,10$. The
script \texttt{paperC/verify\_sticky\_fisher.py} computes the three moments exactly for
$N\le12$ by a dynamic program with $q$ as a symbol, and every difference is
$0$. For fixed $N$ both sides
are rational in $q$ with no pole at $\frac12$, so the case $q=\frac12$
follows. Finally, with $R$ standing for $\rho^N$ and $\rho R$ for
$\rho^{N+1}$, the claim $\E X^4-(\E X^2)^2=2N^2p^2V/q^2$ is an identity of
rational functions of $N$, $q$ and $R$, which the same script checks with
sympy.
\end{proof}

\begin{proof}[Proof of the closed form in Theorem~\ref{thm:fisher}]
For $N\ge2$, $X^2$ is $N^2$ on the constant paths and $(N-2)^2$ on the
paths that switch only at the last step, so $\Var X^2>0$ and $V>0$ by
Lemma~\ref{lem:fisher-moments}, which also gives
$\Cov(J,X^2)=-\frac{Np}q[1+\rho^N-p(1-\rho^N)/T]$, so
$\Cov(J,X^2)^2/\Var X^2$ is the squared bracket over $2V$. We divide by
$\Var J=(N-1)pq$.
\end{proof}

\begin{lemma}[convexity]\label{lem:fisher-convex}
Let $f$ and $g$ be convex and differentiable on $[t-\delta,t+\delta]$, with
$0\le g-f\le\eta$ and $g''\le\kappa$ there. Then
$|f'(t)-g'(t)|\le\eta/\delta+\kappa\delta/2$.
\end{lemma}

\begin{proof}
By convexity and Taylor's theorem,
$f'(t)\le(f(t+\delta)-f(t))/\delta\le(g(t+\delta)-g(t)+\eta)/\delta
\le g'(t)+\kappa\delta/2+\eta/\delta$, and the lower bound is the same with
$t-\delta$.
\end{proof}

\begin{lemma}[a Bessel ratio]\label{lem:fisher-bessel}
For $t>0$ let $R(t)=I_1(t)/I_0(t)$. Then
$(\sqrt{1+t^2}-1)/t\le R(t)<1$. In particular $0<1-R(t)\le1/t$.
\end{lemma}

\begin{proof}
Since $I_0'=I_1$ and $I_1'(t)=I_0(t)-I_1(t)/t$, we have $R'=1-R/t-R^2$, and
$R(t)=t/2+O(t^3)$ at $0$. At a first point $t_0$ with $R(t_0)=1$ we would
have $R'(t_0)\ge0$, but the equation gives $-1/t_0$. So $R<1$. Put
$s=\sqrt{1+t^2}$ and $\ell=(s-1)/t$. Then $1-\ell/t-\ell^2=\ell/t=1/(s+1)$ and
$\ell'=1/(s(s+1))$, so $\ell'\le1-\ell/t-\ell^2$. For $d=R-\ell$ this gives
$d'\ge-d/t-(R^2-\ell^2)=-d(1/t+R+\ell)$. So $\Theta d$ is nondecreasing, where
$\Theta(t)=t\exp\int_0^t(R+\ell)$. As $\ell=t/2+O(t^3)$, we have
$d=O(t^3)$ and $\Theta d\to0$ at $0$. Hence $d\ge0$. Finally
$1-\ell=(t+1-s)/t\le1/t$.
\end{proof}

\begin{lemma}[moments of the endpoint]\label{lem:fisher-W}
\begin{enumerate}[(a)]
\item For each $m\ge1$ there is $C_m$ with $\E W^{2m}\le C_mT^{-m}$ for all
$0<q<1$ and $T\ge1$.
\item If $0<q\le\frac12$, then $\E W^4\le3p^2/T^2+9/(2T^4)$.
\end{enumerate}
\end{lemma}

\begin{proof}
(a) Add a step $\sigma_{N+1}$ to the walk. Since
$\E[\sigma_{i+1}\mid\sigma_{\le i}]=\rho\sigma_i$, the
$\xi_{i+1}=\sigma_{i+1}-\rho\sigma_i$ are martingale differences. Summing
$\sigma_{i+1}-\sigma_i=-2q\sigma_i+\xi_{i+1}$ over $i\le N$ gives
$X=(\sigma_1-\sigma_{N+1}+\Xi)/(2q)$ with $\Xi=\sum_{i=2}^{N+1}\xi_i$. Also
$\xi_{i+1}^2$ is $4q^2$ if step $i+1$ repeats step $i$ and $4p^2\le4$
otherwise. So $\sum\xi_i^2\le4q^2N+4J'$, where $J'\sim\operatorname{Bin}(N,q)$
counts the switches among the $N+1$ steps. Burkholder's
inequality~\cite[Ch.~2]{hallheyde1980} gives
$\E\,\Xi^{2m}\le c_m\E(4q^2N+4J')^m\le c_m8^m((q^2N)^m+\E J'^m)$, and
$q^2N\le T$. The falling factorial moments of $J'$ are at most $T^j$, and
$J'^m$ is a combination of the falling factorials of order $j\le m$ with
nonnegative coefficients. So
$\E J'^m\le b_mT^m$ for $T\ge1$. As $|\sigma_1-\sigma_{N+1}|\le2$,
we get $\E X^{2m}\le(2q)^{-2m}2^{2m-1}(2^{2m}+\E\,\Xi^{2m})\le C_mq^{-2m}T^m$,
and $\E W^{2m}=N^{-2m}\E X^{2m}\le C_mT^{-m}$.

(b) We divide the formula for $\E X^4$ by $N^4$. As $0\le\rho<1$, the two terms
of order $T^{-3}$ are at most $0$ and the one of order $T^{-4}$ is at most
$9/(2T^4)$.
\end{proof}

\begin{proof}[Proof of the bound in Theorem~\ref{thm:fisher}]
Let $\mathcal P$ be the orthogonal projection in $L^2$ onto the span of $1$
and $X^2$. Since $X^2$ is a function of $K$, $\Cov(Z,X^2)=\Cov(J,X^2)$. So
$\Var Z=\Cov(J,X^2)^2/\Var X^2+\E(Z-\mathcal PZ)^2$, and
$e_N-R^2_N=\E(Z-\mathcal PZ)^2/\Var J\ge0$. Also
$\E(Z-\mathcal PZ)^2\le\E(Z-G)^2$ for $G=y(1-W^2/2)+\frac12$. From now on
$0<q\le\frac12$ and $1\le T\le N^{1/4}$, and $C$ is an absolute constant that
may change at each use. Then $p\ge\frac12$, $T\le y\le2T$ and
$\Var J=(N-1)pq\ge T/4$. It suffices to show
$\E(Z-G)^2\le C(T^{-2}+T^4/N)$.

The plan is as follows. By Proposition~\ref{prop:eta}, $Z$ is the logarithmic
derivative of the polynomial $S_{k,N-k}$ at $u$. Paper~A gives a Bessel
function that is close to this polynomial. Closeness of two functions does not
give closeness of their derivatives in general, but it does for convex
functions, and this is Step~1. Step~2 computes the logarithmic derivative of
the Bessel function and finds $y\sqrt{1-W^2}+\frac12$ up to a small error.
Step~3 replaces the square root by $1-W^2/2$. Steps~4 and~5 add up the errors
in mean square.

For $0<k<N$ put $s_k=\sqrt{k(N-k)}$, $c_k=N/(2s_k)$ and $w_k=2s_ku$. On
$K=k$ we have $s_k^2=N^2(1-W^2)/4$, so $c_k=(1-W^2)^{-1/2}\ge1$,
$w_k=y\sqrt{1-W^2}$ and $c_kw_k=y$. Let $\mathcal I_k(z)=I_0(2z)+c_kI_1(2z)$
and $\widetilde S_k(v)=2v\,\mathcal I_k(s_kv)$, which approximates
$S_{k,N-k}(v)$. Here $I_0$ and $I_1$ are the modified Bessel functions. By
the uniform Bessel estimate for every bin in~\cite{paperA} (the theorem
``uniform estimate for every bin'' there, with $a=k$ and $b=N-k$),
\begin{equation}\label{eq:fisher-bessel-bound}
 0\le1-\frac{S_{k,N-k}(v)}{\widetilde S_k(v)}\le\frac{Nv^2}2+v
 \qquad\text{for all }v>0 .
\end{equation}
Let $\Lambda_k(t)=\log S_{k,N-k}(e^t)$ and
$\beta_k(t)=\log\widetilde S_k(e^t)$. By Proposition~\ref{prop:eta},
$Z=\Lambda_K'(\log u)$ on $0<K<N$.

\emph{Step 1.} Both $\Lambda_k$ and $\beta_k$ are convex, as logarithms of
series in $e^t$ with nonnegative coefficients. Let $|t-\log
u|\le1$. Then $v=e^t\le eu$ and $y\le2N^{1/4}$ give
$Nv^2/2+v\le(e^2y^2/2+ey)/N\le(2e^2+2e)/\sqrt N\le\frac12$ for $N\ge1700$. As
$-\log(1-\varepsilon)\le2\varepsilon$ for $0\le\varepsilon\le\frac12$,
\eqref{eq:fisher-bessel-bound} gives $0\le\beta_k-\Lambda_k\le\eta$ there,
with $\eta=(e^2y^2+2ey)/N$. Next, $\beta_k(t)-t-\log2$ is the logarithm of
$\mathcal I_k(s_ke^t)=\sum_na_ne^{nt}$ with $a_n\ge0$. So $\beta_k''$ is the
variance of $n$ under the weights $a_ne^{nt}$, at most the mean of $n^2$. As
$(z\frac d{dz})^2$ maps $I_0(2z)$ to $4z^2I_0(2z)$ and $I_1(2z)$ to
$(4z^2+1)I_1(2z)$, this mean is at most $4z^2+1\le e^2y^2+1=\kappa$, since
$z=s_ke^t\le ey/2$. For $N\ge4$,
$\delta=\sqrt{2\eta/\kappa}\le1$ and Lemma~\ref{lem:fisher-convex} gives
\begin{equation}\label{eq:fisher-step1}
 \bigl|\Lambda_k'(\log u)-\beta_k'(\log u)\bigr|
 \le\sqrt{2\eta\kappa}
 \le\frac{2e^2y^2+2ey+1}{\sqrt{2N}}\le\frac{11(1+y)^2}{\sqrt N}.
\end{equation}

\emph{Step 2.} Write $z=s_ke^t$, $w=2z$ and $R=R(w)$. By the formulas for
$I_0'$ and $I_1'$,
\[
 \beta_k'(t)=1+\frac{z\mathcal I_k'(z)}{\mathcal I_k(z)}
 =1+\frac{wR+c_kw-c_kR}{1+c_kR}=w+\frac{1+w(c_k-1)(1-R)}{1+c_kR}.
\]
At $t=\log u$ we have $w=w_k$, so $\beta_k'(\log u)=w_k+B_k$ with $B_k$ the
last fraction. By Lemma~\ref{lem:fisher-bessel}, $0<R<1$ and $w(1-R)\le1$.
So the numerator of $B_k$ lies in $[1,c_k]$ and is at most
$1+wc_k=1+y$. Hence $0<B_k\le1+y$. Subtracting $\frac12$ from
$1/(1+c_kR)\le B_k\le c_k/(1+c_kR)$ gives
\[
 -\frac{c_k-1}2\le B_k-\frac12\le\frac{(c_k-1)+c_k(1-R)}{2(1+c_kR)}
 \le\frac{c_k-1}2+\frac{c_k^2}{2y},
\]
by $c_k(1-R)\le c_k/w_k=c_k^2/y$. If $W^2\le\frac12$ on $K=k$, then
$c_k^2\le2$ and $c_k-1\le W^2$, since $(1-v)^{-1/2}-1$ is convex, vanishes
at $0$ and is below $\frac12$ at $v=\frac12$. So
$|B_k-\frac12|\le W^2/2+1/y$ there.

\emph{Step 3.} Let $\psi(a)=1-a^2/2-\sqrt{1-a^2}$ for $|a|\le1$.
The series of $\sqrt{1-v}$ is $1-v/2-\sum_{j\ge2}b_jv^j$ with $b_j\ge0$ and
$\sum_jb_j=\frac12$, so $0\le\psi(a)\le a^4/2$. On $K=k$,
$w_k=y(1-W^2/2)-y\psi(W)$. So on $0<K<N$ we have
$Z-G=(Z-\beta_K'(\log u))+(B_K-\frac12)-y\psi(W)$, and on the atoms
$K\in\{0,N\}$ we have $Z=0$, $W^2=1$ and $Z-G=-(y+1)/2$.

\emph{Step 4.} By Lemma~\ref{lem:fisher-W} and Markov's inequality,
$\E W^4\le8/T^2$, $\E W^8\le CT^{-4}$ and
$\Prob(W^2>\frac12)\le16\,\E W^8\le CT^{-4}$. Also
$\Prob(K\in\{0,N\})=p^{N-1}\le e^{-q(N-1)}\le e^{1/2-T}$.

\emph{Step 5.} By Step~3 and Minkowski's inequality, $\lVert Z-G\rVert_2$ is
at most the sum of the $L^2$ norms of the three terms of Step~3 on $0<K<N$ and of
$\frac{y+1}2$ on the atoms. By \eqref{eq:fisher-step1} the first is at most
$99T^2/\sqrt N$, as $(1+y)^2\le9T^2$. For the second we use Step~2 where
$W^2\le\frac12$ and $|B_K-\frac12|\le1+y$ elsewhere, so its square is at most
$\E W^4/2+2/y^2+C(1+y)^2T^{-4}\le C/T^2$. The third is at most
$\frac y2(\E W^8)^{1/2}\le C/T$, and the last is
$\frac{y+1}2\Prob(K\in\{0,N\})^{1/2}\le CTe^{-T/2}\le C/T$. So
$\E(Z-G)^2\le C(T^4/N+T^{-2})$, and dividing by $\Var J\ge T/4$ gives the
bound. Only Step~1 used $N\ge1700$, so
$N_0=1700$ works.
\end{proof}


\begin{thebibliography}{9}\small

\bibitem{kagey131}
P.~Kagey, \emph{Problem 131}, Open Problem Collection,
\url{https://peterkagey.com/problems/131/} (accessed October 2026).

\bibitem{paperA}
A.~Pemmasani, \emph{Uniform Bessel bounds for endpoint crossings in the
persistent random walk}, preprint, 2026.
\href{https://arxiv.org/abs/2610.06935}{arXiv:2610.06935}.

\bibitem{hallheyde1980}
P.~Hall and C.~C.~Heyde, \emph{Martingale Limit Theory and Its
Application}, Academic Press, New York, 1980.

\bibitem{smiga2023}
J.~A.~Smiga, M.~Radaelli, F.~C.~Binder and G.~T.~Landi, \emph{Stochastic
metrology and the empirical distribution}, Phys.~Rev.~Research \textbf{5}
(2023), 033150. \url{https://doi.org/10.1103/PhysRevResearch.5.033150}.

\end{thebibliography}

\end{document}
